% Template for Elsevier CRC journal article
% version 1.1-5p dated 18 January 2011
% SEM ACENTUACAO PROBLEMA UTF-8 : NAO ACEITA

% This file (c) 2010-2011 Elsevier Ltd.  Modifications may be freely made,
% provided the edited file is saved under a different name

% This file contains modifications for Procedia Computer Science
% but may easily be adapted to other journals

% Changes since version 1.0
% - elsarticle class option changed from 1p to 3p (to better reflect CRC layout)
% - this version uses option 5p for larger-format journals (text area 24.1 x 18.4 cm)

%-----------------------------------------------------------------------------------

%% This template uses the elsarticle.cls document class and the extension package ecrc.sty
%% For full documentation on usage of elsarticle.cls, consult the documentation "elsdoc.pdf"
%% Further resources available at http://www.elsevier.com/latex

%-----------------------------------------------------------------------------------

%%%%%%%%%%%%%%%%%%%%%%%%%%%%%%%%%%%%%%%%%%%%%%
%%%%%%%%%%%%%%%%%%%%%%%%%%%%%%%%%%%%%%%%%%%%%%
%%                                          %%
%% Important note on usage                  %%
%% -----------------------                  %%
%% This file must be compiled with PDFLaTeX %%
%% Using standard LaTeX will not work!      %%
%%                                          %%
%%%%%%%%%%%%%%%%%%%%%%%%%%%%%%%%%%%%%%%%%%%%%%
%%%%%%%%%%%%%%%%%%%%%%%%%%%%%%%%%%%%%%%%%%%%%%

%% The '5p' and 'times' class options of elsarticle are used for Elsevier CRC
\documentclass[final,5p,times,twocolumn,authoryear,10pt]{elsarticle}

%% The `ecrc_jnr' package must be called to make the CRC functionality available according to the JNR style.
%% ecrc_jnr es el paquete ecrc de Elsevier con modificaciones para JNR2018
\usepackage{ecrc_jnr}

%% The ecrc package defines commands needed for running heads and logos.
%% For running heads, you can set the journal name, the volume, the starting page and the authors

%%%%%%%%%%%%%%%%%%%%%%%%%%%%%%%%% Añadido por Secretaría RIAI
\usepackage[spanish]{babel}     % Idioma
\addto\captionsspanish{%
\def\tablename{Tabla}%
}
\usepackage[utf8]{inputenc}   % Lengua Castellano
\usepackage{amsmath}            % Para las referencias a ecuaciones con \eqref
\usepackage{epstopdf}           % Para poder insertar figuras .eps al compilar con PDFLATEX
\usepackage{flushend}           % Para igualar las columnas de la última página
%\usepackage{hyperref}           % Para hipervínculos dentro del PDF

\usepackage{graphicx}
\graphicspath{{./Imagenes/}}
\DeclareGraphicsExtensions{.jpeg,.png}

\usepackage{array}

%%%%%%%%%%%%%%%%%%%%%%%%%%%%%%%%%%%%%%%%%%%%%%%%%%%%%%%


%% set the volume if you know. Otherwise `00'
\volume{2018}

%% set the starting page if not 1
\firstpage{1}

%% Give the name of the journal
\journalname{Jornadas Nacionales de Robótica}

%% Give the author list to appear in the running head
%% Example \runauth{C.V. Radhakrishnan et al.}
%\runauth{Primer autor et al.}

%% The choice of journal logo is determined by the \jid and \jnltitlelogo commands.
%% A user-supplied logo with the name <\jid>logo.pdf will be inserted if present.
%% e.g. if \jid{yspmi} the system will look for a file yspmilogo.pdf
%% Otherwise the content of \jnltitlelogo will be set between horizontal lines as a default logo

%% Give the abbreviation of the Journal. Contast the Publisher if in doubt what this is.
\jid{JNR}

%% Give a short journal name for the dummy logo (if needed)
\jnltitlelogo{}

%% End of ecrc-specific commands
%%%%%%%%%%%%%%%%%%%%%%%%%%%%%%%%%%%%%%%%%%%%%%%%%%%%%%%%%%%%%%%%%%%%%%%%%%

%% The amssymb package provides various useful mathematical symbols
\usepackage{amssymb}
%% The amsthm package provides extended theorem environments
%% \usepackage{amsthm}

%% The lineno packages adds line numbers. Start line numbering with
%% \begin{linenumbers}, end it with \end{linenumbers}. Or switch it on
%% for the whole article with \linenumbers after \end{frontmatter}.
%% \usepackage{lineno}

%% natbib.sty is loaded by default. However, natbib options can be
%% provided with \biboptions{...} command. Following options are
%% valid:

%%   round  -  round parentheses are used (default)
%%   square -  square brackets are used   [option]
%%   curly  -  curly braces are used      {option}
%%   angle  -  angle brackets are used    <option>
%%   semicolon  -  multiple citations separated by semi-colon
%%   colon  - same as semicolon, an earlier confusion
%%   comma  -  separated by comma
%%   numbers-  selects numerical citations
%%   super  -  numerical citations as superscripts
%%   sort   -  sorts multiple citations according to order in ref. list
%%   sort&compress   -  like sort, but also compresses numerical citations
%%   compress - compresses without sorting
%%
%% \biboptions{comma,round}

% \biboptions{}

% if you have landscape tables
\usepackage[figuresright]{rotating}

% put your own definitions here:
%   \newcommand{\cZ}{\cal{Z}}
%   \newtheorem{def}{Definition}[section]
%   ...

% add words to TeX's hyphenation exception list
%\hyphenation{author another created financial paper re-commend-ed Post-Script}

% para poder introducir varias figuras que ocupen el ancho de las dos columnas.
\usepackage{subfigure}

\usepackage{lipsum}
\usepackage{setspace}

% declarations for front matter
\begin{document}
\fontsize{9.5pt}{11.4pt}\selectfont
\begin{frontmatter}

%% Title, authors and addresses

%% use the tnoteref command within \title for footnotes;
%% use the tnotetext command for the associated footnote;

%% use the fnref command within \author or \address for footnotes;
%% use the fntext command for the associated footnote;

%% use the corref command within \author for corresponding author footnotes;
%% use the cortext command for the associated footnote;
%% use the ead command for the email address,
%% and the form \ead[url] for the home page:
%%
%% \tnotetext[label1]{}
%% \author{Name\corref{cor1}\fnref{label2}}
%% \ead{email address}
%% \ead[url]{home page}
%% \fntext[label2]{}
%% \cortext[cor1]{}
%% \address{Address\fnref{label3}}
%% \fntext[label3]{}

%\dochead{Cabecera artículo}
%% Use \dochead if there is an article header, e.g. \dochead{Short communication}

\title{
\fontsize{15pt}{18pt}\selectfont
%título de la comunicación
Preparación de Artículos para la revista RIAI:\\
Use Tipo Título para el Título del Artículo}


%% use optional labels to link authors explicitly to addresses:
%% \author[label1,label2]{<author name>}
%% \address[label1]{<address>}
%% \address[label2]{<address>}

\author[First]{Primer A. Autor\corref{cor1}\fnref{label2}}
\ead{autor@cea-ifac.es}
\ead[url]{www.cea-ifac.es}

\author[Second]{Segundo B. Autor, Jr.}
\ead{autor2@cea-ifac.es}
\ead[url]{www2.cea-ifac.es}

\author[Third]{Tercer C. Autor}
\ead{autor3@cea-ifac.es}

\cortext[cor1]{Autor en correspondencia.}
\fntext[label2]{Nota al pie para el autor 1}



\address[First]{
\fontsize{7.5pt}{11.4pt}\selectfont
Comité Español de Automática, Parc Tecnologic de Barcelona, Edifici U, C/ Llorens i Artigas, 4-6, 08028 Barcelona, España}
\address[Second]{
\setstretch{0.4}
\fontsize{7.5pt}{11.4pt}\selectfont
Departamento de Automática, Ingeniería Electrónica e Informática,Universidad Politécnica de Madrid, C/ José Gutiérrez Abascal, nº2, 28006, Madrid, España}
\address[Third]{
\setstretch{0.8}
\fontsize{7.5pt}{11.4pt}\selectfont
Robotics Lab, Departmento de Ingeniería de Sistemas y Automática, Universidad Carlos III de Madrid, Avda. de la Universidad 30, 28911 Leganés, España}


\begin{abstract}
\setstretch{0.96}
\fontsize{9.5pt}{11.4pt}\selectfont
%% Text of abstract
\noindent\hspace{0.3cm}
Estas instrucciones constituyen una guia para la preparación de artículos para la revista RIAI. Utilice este documento como un conjunto de instrucciones. Por favor, utilice este documento como una "plantilla" para preparar su manuscrito. Para las directrices de envío, siga las instrucciones del sistema de envío de artículos de la página web de la revista. \emph{Copyright {\copyright} XXXX CEA.}
\end{abstract}

\begin{keyword}
%% MSC codes here, in the form: \MSC code \sep code
%% or \MSC[2008] code \sep code (2000 is the default)
\setstretch{0.96}
\fontsize{9.5pt}{11.4pt}\selectfont

%% keywords here, in the form: keyword \sep keyword
palabra 1 \sep palabra 2 \sep 5-10 palabras clave (tomadas de la lista del sitio web de IFAC).

%******************************************************************************
% Comentar Sección siguiente si la comunicación es de proyecto o tesis doctoral
%******************************************************************************
\textit{Datos del Proyecto:}

\hspace{1.26cm}\textit{Denominaci\'on del proyecto:} Nombre del proyecto...

\hspace{1.26cm}\textit{Referencia:}

\hspace{1.26cm}\textit{Investigador/es responsable/es:}

\hspace{1.26cm}\textit{Tipo de proyecto (internacional, nacional, auton\'omico, transferencia):}

\hspace{1.26cm}\textit{Entidad/es financiadora/s:}

\hspace{1.26cm}\textit{Fecha de inicio/fin:}

%***********************************************************
% Para el caso de tesis doctoral
Alternativamente:

\textit{Datos de la tesis:} Título de la tesis...

\hspace{1.26cm}\textit{Doctorando:}

\hspace{1.26cm}\textit{Directores:}

\hspace{1.26cm}\textit{Fecha Inicio/Fecha estimada terminaci\'on:}

\hspace{1.26cm}\textit{Programa de doctorado:}

\hspace{1.26cm}\textit{Universidad:}
%************************************************************

\end{keyword}

\end{frontmatter}
%%
%% Start line numbering here if you want
%%
% \linenumbers
\setlength{\parindent}{0.42cm}
%% main text
\section{Introducción}
Estas instrucciones constituyen una guía para la preparación de artículos para la revista RIAI. Utilice este documento como un conjunto de instrucciones. Por favor, utilice este documento como una “plantilla” para preparar su manuscrito.

Para las directrices de envío, siga las instrucciones del sistema de envíode artículos de la página web de la revista. \textbf{No cambie el tamaño de las fuentes o espaciado de línea para introducir más texto en un número limitado de páginas.} Utilice cursiva para enfatizar; no subraye. 

\subsection{Una subsección}

Bifurcación: Trazado del máximo local de x con una disminución de amortiguamiento $a$ (Figura \ref{fig1}).

Para insertar imágenes en \textit{Word}, posicione el cursor en el punto de inserción y o bien utilice Insertar - Imagen -  Desde Fichero o copie la imagen al portapapeles de Windows y entonces use Editar - Pegado especial - Imagen (con “Flotar sobre el texto” deseleccionado).

\begin{figure}
\centering
  % El fichero es un eps y se convierte automaticamente a pdf con eps2pdf package
  \includegraphics[width=7cm]{figura1.jpg}\\
  \caption{Bifurcación. Trazado del máximo local de $x$ con una disminución de amortiguamiento $a$.}
  \label{fig1}
\end{figure}

RIAI no realizará ninguna operación de formateado final  a su artículo. Su documento debe estar “listo para filmar”. El número límite de hojas para el documento es de doce. \textbf{Por favor, no modifique los márgenes. Si está creando el documento usted mismo, tenga en cuenta los márgenes listados en la Tabla 1.}

\section{Procedimiento para el Envío de Artículos}
Recuerde que RIAI está considerado como un “Cámera Ready Copy Journal” (CRC). Esto implica que los autores son responsables de aplicar el formato correspondiente a sus contribuciones. Desde la secretaría de la revista no se ejecutará ninguna acción de formateo a los artículos. A continuación vemos unas subsecciones.

\subsection{Fase de Revisión}
Por favor, use este documento como una “plantilla” para preparar su documento. Para las directrices de envío, siga las instrucciones del sistema de envío de artículos. 

Dado que el límite de páginas es de seis, es mejor preparar el envío inicial en el formato listo para filmar, de tal manera que tenga una buena estimación de la longitud de hojas. Adicionalmente, el esfuerzo requerido para el envío de la versión final será, de esta manera, mínimo.

\begin{figure}
\centering
  % Se pueden incluir figuras jpeg al compilar con PDFLatex
  \includegraphics[width=3cm]{figurajpeg.jpg}\\
  \caption{Título de la figura 2}\label{fig2}
\end{figure}


\subsection{Fase Final}
Se supone que los autores tendrán en cuenta rigurosamente los márgenes. En caso de no ser así se le pedirá que reenvíe el documento para que así lo cumpla, retrasando de esta manera la preparación de los contenidos de la revista. 

\subsection{Márgenes de página}
Todas las dimensiones están en centímetros.

\begin{table}[htbp]
  \caption{Márgenes de página (cm.)}
   \label{tabla1}
   \centering
  \begin{tabular}{cccc}
\hline
Página &	Arriba	& Abajo & Izq./Der. \\
    \hline
Primera	& 4 & 2.5 &	 1.5 \\
Resto	& 2.5 & 2.5 &  1.5 \\
\hline
  \end{tabular}
\end{table}

Es muy importante mantener estos márgenes. Son necesarios para poner información de la revista y los números de página.

\subsection{Figuras y Creación del  PDF}
Todas las figuras deben estar incrustadas en el documento. Cuando incluya una imagen, asegúrese de insertar la imagen real en lugar de un enlace a su computador local. En la medida de lo posible, utilice las herramientas de conversión a PDF estándares Adobe Acrobat o Ghostscript que dan los mejores resultados. \textbf{Es importante que todas las fuentes estén incrustadas/subconjunto en el PDF resultante.}

\section{Matemáticas}
Si utiliza Word, use o bien Microsoft Editor de Ecuaciones o  MathType  para las ecuaciones de su artículo (Insertar - Objeto - Crear Nuevo - Microsoft Editor de Ecuaciones o Ecuación MathType). No debe seleccionarse la opción “Flotar” sobre el texto. Por supuesto, LaTeX gestiona las ecuaciones a través de macros pre-programadas.

\section{Unidades}

Use el Sistema Internacional como unidades primarias. Se pueden usar otras unidades como unidades secundarias (entre paréntesis). Esto se aplica a artículos sobre almacenamiento de datos. Por ejemplo, escriba “15 Gb/cm2 (100 Gb/in2).” Se considera una excepción cuando las unidades inglesas se usan como identificadores comerciales, como “unidad de disco de 3.5 pulgadas.” Evite mezclar unidades del Sistema Internacional con el Sistema Cegesimal, tales como corriente en amperios y campo magnético en  oersteds. Esto a menudo lleva a confusión porque las ecuaciones no son dimensionalmente equiparables. Si debe usar unidades mezcladas, especifique claramente las unidades para cada cantidad  en la ecuación.

La unidad en el Sistema Internacional para la fuerza del campo magnético H es A/m. Sin embargo, si desea utilizar unidades de $T$, o bien refiérase a densidad de flujo magnético B o fuerza del campo magnético simbolizado como $B$ o fuerza del campo magnatico simbolizado como $\mu_0 H$. Utilice el punto centrado para separar unidades compuestas, es decir, $A\cdot m^2$.

\section{Consejos Útiles}
\subsection{Figuras y Tablas}
Las etiquetas de los ejes de las figures son a menudo fuentes de confusión. Utilice palabras en lugar de símbolos. Como ejemplo, escriba la cantidad “Magnetización,” o “Magnetización M,” no sólo “M.” Ponga las unidades entre paréntesis. No etiquete los ejes únicamente con unidades. Como en la Figura 1, por ejemplo, escriba “Magnetización (A/m)” o “Magnetización (A m1),” no sólo “A/m.” No etiquete los ejes con una relación de cantidades y unidades. Por ejemplo, escriba “Temperatura (K),” no “Temperatura/K.”

Los multiplicadores pueden ser especialmente Fuente de confusión. Escriba “Magnetización (kA/m)” o “Magnetización (103 A/m).” No escriba “Magnetización (A/m)  1000” porque el lector no sabría si la etiqueta del eje superior en la Fig. 1 es 16000 A/m o 0.016 A/m. Las etiquetas de las figuras deben ser legibles, aproximadamente de 8 a 12 puntos.

\subsection{Referencias}
La lista de referencias debe ser ordenada alfabéticamente de acuerdo con el primer autor, con las siguientes líneas justificadas con la sangría correspondiente. Si existen diferentes publicaciones del mismo autor(es), éstas deberán ser listadas en el orden del año de publicación. Si hay más de un artículo del mismo autor en la misma fecha, etiquételas como a,b, etc. (Baker, 1990a, b). Por favor, fíjese que todas las referencias (García, 2007) en este apartado (García and Martínez, 2008) deben ser citadas directamente en el cuerpo del texto (García et al., 2004). Algunos ejemplos: \cite{Bak:63a,Bak:63b}, \cite{Keohane:58},\cite{ChaRou:66}.

Por favor, tenga en cuenta que las referencias al final de este documento cumplen con el estilo anteriormente mencionado. Los artículos que no hayan sido publicados deben ser citados como “no publicado.” Ponga en mayúscula únicamente la primera palabra del título, excepto el caso de nombres propios y símbolos de elementos.

Si está utilizando LaTeX, puede procesar una base de datos de bibliografía externa o insertarla directamente en la sección de referencias. Las notas al pie de página se deben evitar en la medida de lo posible.

\subsection{Abreviaciones y Acrónimos}
Defina las abreviaciones y acrónimos la primera vez que se usan en el texto, incluso después de que ya hayan sido definidos en el resumen. Abreviaciones tales como IFAC, SI, ac, y dc no necesitan ser definidas. Abreviaciones que incorporen periodos no deben tener espacios: escriba “C.N.R.S.,” no “C. N. R. S.” No utilice abreviaciones en el título salvo que sea inevitable (por ejemplo, “RIAI” en el título de este artículo).


\subsection{Ecuaciones}

Numere las ecuaciones consecutivamente con números de ecuaciones entre paréntesis justificado al margen derecho, como en \eqref{e2}. Primero use el editor de ecuaciones para crear la ecuación. Después seleccione el estilo “Equation”. Presione la tecla de tabulador y escriba el número de ecuación entre paréntesis. Para hacer sus ecuaciones más compactas, puede usar el solidus ( / ), la función exp, o los exponentes apropiados. Utilice paréntesis para evitar ambigüedades en los denominadores. Ponga signos de puntuación en las ecuaciones cuando formen parte de una frase, como en

\begin{equation} \label{e2}
\begin{array}{ll}
\int_0^{r_2} & F (r, \varphi ) dr d\varphi = [\sigma r_2 / (2 \mu_0 )] \\
& \cdot \int_0^{\inf} exp(-\lambda |z_j - z_i |) \lambda^{-1} J_1 (\lambda  r_2 ) J_0 (\lambda r_i ) d\lambda
\end{array}
\end{equation}

Asegúrese de que los símbolos de su ecuación han sido definidos antes de que la ecuación aparezca o inmediatamente después. Ponga en cursiva los símbolos (T podría referirse a la temperatura, pero T es la unidad tesla). Refiérase a “(1),” no “Ec. (1)” o “ecuación (1),” excepto al principio de la frase: “La ecuación (1) es ... .”

\subsection{Otras Recomendaciones}

Utilice un espacio tras los periodos y dos puntos. Evite utilizar participios, tales como, “Utilizando (1), se calculó el potencial.” [No está claro quién o qué usó (1).] En su lugar escriba, “El Potencial fue calculado empleando (1),” o “Empleando (1) se calculó el potencial.”

\section{Conclusión}

Una sección de conclusiones no es necesaria. Sin embargo, las conclusiones pueden revisar los puntos más importantes de un artículo, pero no debe replicarse el resumen en las conclusiones. Las conclusiones pueden tratar sobre la importancia del trabajo realizado o sugerir aplicaciones o trabajos futuros.
%%%%%%%%%%%%%%%%%%%%%%%%%%%%%%%%%%%%%%%%%%%%%%%%%%%%%%%%%%%%%%%%%%%%%%%%%%%%%%%%%%%%%%%%
%% ATENCION AUTORES: ESTA SECCION ES OBLIGATORIA
%%%%%%%%%%%%%%%%%%%%%%%%%%%%%%%%%%%%%%%%%%%%%%%%%%%%%%%%%%%%%%%%%%%%%%%%%%%%%%%%%%%%%%%%

\section*{English Summary}

\textbf{Paper title in English, bold style.}\\

\noindent \textbf{Abstract}\\

Physical Programming (PP) is a multiobjective optimization technique where the designer, for each objective or specification of the problem, declares his preferences in a flexible way. These preferences(e.g. overshoot, settling time, gain margin) are expressed with linguistic terms such as tolerable, desirable, undesirable, etc. and they are associated with numeric ranges in the same physical units as the objective itself (e.g. seconds, percentages). This paper shows how PP can be applied to the design of robust controllers from a multiobjective optimization point of view. Non linear optimization used in the original PP method has been substituted by a Genetic Algorithm to avoid local minima which can usually arise in this multimodal problems. The ACC Robust Control Benchmark has been solved and the result obtained is compared with solutions from other authors.\\

\noindent \emph{Keywords:}\\

Robust control, multiobjective optimization, Physical Programming, Genetic Algorithms.
%%%%%%%%%%%%%%%%%%%%%%%%%%%%%%%%%%%%%%%%%%%%%%%%%%%%%%%%%%%
\section*{Agradecimientos}

Este trabajo ha sido realizado parcialmente gracias al apoyo de la Agencia Nacional (los agradecimientos de financiación y apoyos han de ser incluidos aquí).

\label{}

%% The Appendices part is started with the command \appendix;
%% appendix sections are then done as normal sections
%% \appendix

%% \section{}
%% \label{}

%% References
%%
%% Following citation commands can be used in the body text:
%% Usage of \cite is as follows:
%%   \cite{key}         ==>>  [#]
%%   \cite[chap. 2]{key} ==>> [#, chap. 2]
%%


%% References with BibTeX database:

\bibliographystyle{elsarticle-harv}
%%\bibliography{riaibib}

%% Authors are advised to use a BibTeX database file for their reference list.
%% The provided style file elsarticle-harv.bst formats references in the required Procedia style

%% For references without a BibTeX database:

%%\begin{thebibliography}{01}

%% \bibitem must have the following form:
%%   \bibitem{key}...
%%

% \bibitem{}

%%\bibitem[{Able(1945)}]{Abl:45}
%%Able, B., 1945. Nombre del artículo. Nombre de la revista 35, 123--126.

\begin{thebibliography}{7}
\expandafter\ifx\csname natexlab\endcsname\relax\def\natexlab#1{#1}\fi
\expandafter\ifx\csname url\endcsname\relax
  \def\url#1{\texttt{#1}}\fi
\expandafter\ifx\csname doi\endcsname\relax
  \def\doi#1{\texttt{#1}}\fi
\expandafter\ifx\csname urlprefix\endcsname\relax\def\urlprefix{URL: }\fi
\expandafter\ifx\csname doiprefix\endcsname\relax\def\doiprefix{DOI: }\fi

\bibitem[{(Able, 1945)}]{Abl:45}
Able, B., 1945. Nombre del artículo. Nombre de la revista 35, 123--126.

\bibitem[{(Able, 1956)}]{Abl:56}
Able, B., 1956. Nombre del artículo. Nombre de la revista 135, 7--9.

\bibitem[{(Baker, 1963{\natexlab{A}})}]{Bak:63a}
Baker, R., 1963{\natexlab{A}}. Nombre de la editorial, Lugar de publicación.

\bibitem[{(Baker, 1963{\natexlab{B}})}]{Bak:63b}
Baker, R., 1963{\natexlab{B}}. Nombre del artículo. Nombre de la revista 34, 184--186.

\bibitem[{(Charlie and Routh, 1966)}]{ChaRou:66}
Charlie, F., Routh, M., 1966. Nombre del artículo. Nombre de la revista 66, 267--269.

\bibitem[{(Dog, 1958)}]{Dog:58}
Dog, P., 1958. In: Editor (Ed.), Título del libro. Vol. II. Nombre de la editorial, Lugar de publicación, Ch. 7, pp. 56–98.

\bibitem[{(García, 2007)}]{Garcia:2007}
García, F., 2007. Nombre del libro. Editorial, Lugar de publicación.

\bibitem[{(García \textit{et al.}, 2004)}]{Garcia:2004}
Garcia, F., López, J., Otro, A., 2004. Nombre del artículo. En: del editor(es), N. (Ed.), Nombre del libro. Editorial, Lugar de publicación, p. números de página.

\bibitem[{(García and Martínez, 2008)}]{Garcia:2008}
García, F., Martínez, R., 2008. Nombre del artículo. Nombre de la revista número, {}números de página.
  
\bibitem[{(Keohane, 1958)}]{Keohane:58}
Keohane, R., 1958. Nombre editorial, Lugar de publicación.

\bibitem[{(Soukhanov, 1992)}]{Heritage:92}
Soukhanov, A.~H. (Ed.), 1992. Nombre de la editorial.

\end{thebibliography}

\appendix
\section{Primer Apéndice}    % Cada Apéndice debe tener un título corto.
Si utiliza Word, use o bien Microsoft Editor de Ecuaciones o  MathType  para las ecuaciones de su artículo (Insertar - Objeto - Crear Nuevo - Microsoft Editor de Ecuaciones o Ecuación MathType). No debe seleccionarse la opción “Flotar” sobre el texto. Por supuesto, LaTeX gestiona las ecuaciones a través de macros pre-programadas.

\section{Segundo Apandice}
Use el Sistema Internacional como unidades primarias. Se pueden usar otras unidades como unidades secundarias (entre paréntesis). Esto se aplica a artículos sobre almacenamiento de datos. Por ejemplo, escriba ``$15 Gb/cm^2$'' ($100 Gb/in^2$). Se considera una excepción cuando las unidades inglesas se usan como identificadores comerciales, como “unidad de disco de 3.5 pulgadas.” Evite mezclar unidades del Sistema Internacional con el Sistema Cegesimal, tales como corriente en amperios y campo magnético en oersteds. Esto a menudo lleva a confusión porque las ecuaciones no son dimensionalmente equiparables. Si debe usar unidades mezcladas, especifique claramente las unidades para cada cantidad en la ecuación.

La unidad en el Sistema Internacional para la fuerza del campo magnético H es A/m. Sin embargo, si desea utilizar unidades de $T$, o bien refiérase a densidad de flujo magnético B o fuerza del campo magnético simbolizado como $B$ o fuerza del campo magnatico simbolizado como $\mu_0 H$. Utilice el punto centrado para separar unidades compuestas, es decir, $A\cdot m^2$.

\section{Tercer Apéndice}

\subsection{Mas sobre Figuras y Tablas}

Las etiquetas de los ejes de las figures son a menudo fuentes de confusión. Utilice palabras en lugar de símbolos. Como ejemplo, escriba la cantidad ``Magnetizacian,'' o ``Magnetización M,'' no sólo ``M.'' Ponga las unidades entre paréntesis. No etiquete los ejes únicamente con unidades. Como en la Figura 1, por ejemplo, escriba ``Magnetización (A/m)'' o ``Magnetización (A $\cdot$ m$^{-1}$),'' no sólo ``A/m'' No etiquete los ejes con una relación de cantidades y unidades. Por ejemplo, escriba ``Temperatura (K),''
no ``Temperatura/K.''

Los multiplicadores pueden ser especialmente Fuente de confusión. Escriba ``Magnetización (kA/m)'' o ``Magnetización ($10^3$ A/m).'' No escriba ``Magnetización (A/m) $\times$ 1000'' porque el lector no sabría si la etiqueta del eje superior en la Fig. 1 es 16000 A/m o 0.016 A/m. Las etiquetas de las figuras deben ser legibles, aproximadamente de 8 a 12 puntos.

\end{document}

%%
%% End of file `ejemplo plantilla latex jnr2018.tex'.